Pour mesurer la distance Terre-Lune, on envoie un faisceau laser de longueur d'onde \qty{\v{lam}}{nm} vers la Lune, à \qty{3.84e8}{m}, à travers un télescope de \qty{\v{Dt}}{m} de diamètre. À cause de la diffraction par l'ouverture du télescope, le faisceau s'élargit : son bord correspond au premier minimum de diffraction d'une ouverture circulaire. On néglige les effets de l'atmosphère.

\begin{enumerate}[label=\alph*)]
    \item Déterminez le demi-angle d'ouverture du faisceau.
    \item Déterminez le diamètre de la tache lumineuse sur la Lune.
    \item Quel diamètre de télescope faudrait-il pour que la tache n'ait que \qty{\v{X}}{km} de diamètre ?
\end{enumerate}
